% Draft for stakeholder validation. This fragment is not yet included in main.tex.
\section{Quantified EPS problem: line of evidence}

\begin{enumerate}
    \item \textbf{External constraint.} ÉTS intends to propose a 2-3U CubeSat under CUBICS 2026 Stream 2. The CSA caps this stream's grant at CAD~\$350{,}000 and its project period at three years, including operations, data collection, and analysis. The grant is a limit on mission funding, not an EPS allocation. Recipients bear any cost overrun. The CSA anticipates launch opportunities three to four years after grant issuance. If the project period starts near grant issuance, the stated three-year duration leaves no funded period for operations after a CSA-procured launch. The call does not explain how these dates fit together. ÉTS must clarify its funding and operations plan with the CSA.\footnote{CSA, \emph{CUBICS 2026 Announcement of Opportunity}, Sections 6.2 and 7.1 and Appendix C.3.2: \url{https://www.asc-csa.gc.ca/eng/funding-programs/funding-opportunities/ao/2026-cubics-cubesats-initiative-in-canada-for-stem.asp}.}

    \item \textbf{Relevant precedent.} The CSA awarded MARMOTSat CAD~\$349{,}600 in the 2022 cohort. Its team reported a tested first full EPS prototype on April 4, 2025, and a tested second revision on November 26, 2025. These reports are 236 days apart. The team reported improved power efficiency in the second revision. The interval establishes that EPS design and validation continued across revisions. It does not measure labor, an avoidable delay, or an overrun.\footnote{CSA, \emph{About CUBICS}: \url{https://www.asc-csa.gc.ca/eng/satellites/cubics/about.asp}. MARMOTSat, \emph{First Full EPS Prototype}: \url{https://www.marmotsat.ca/updates/first-full-eps-prototype}. MARMOTSat, \emph{EPS Second Revision}: \url{https://www.marmotsat.ca/updates/eps-second-revision}.}

    \item \textbf{Process risk.} In its 2026 call, the CSA reports that some teams in earlier CubeSat programs relied heavily on one-year capstone students. Work was sometimes restarted when knowledge was not transferred to the next team. This observation is relevant to ÉTS only if its club has similar staffing and handoff conditions.\footnote{CSA, \emph{CUBICS 2026 Announcement of Opportunity}, Appendix B.3, ``Team Composition'': \url{https://www.asc-csa.gc.ca/eng/funding-programs/funding-opportunities/ao/2026-cubics-cubesats-initiative-in-canada-for-stem.asp}.}

    \item \textbf{Local baseline to establish.} Ask the ÉTS team whether it has previously designed and validated an EPS. If so, obtain the prior spacecraft size (\(X\)U), the EPS scope and test criteria, planned and actual dates, EPS labor hours, EPS purchases and test costs, whole-project costs, and the records behind those figures. Let \(C_{\mathrm{past,EPS}}\) be the actual EPS cost and \(T_{\mathrm{past,EPS}}\) the elapsed time from the defined EPS start to its validation gate. Report the EPS cost share \(A=100C_{\mathrm{past,EPS}}/C_{\mathrm{past,mission}}\) only if whole-project cost is known. Report \(B=100T_{\mathrm{past,EPS}}/T_{\mathrm{past,mission}}\) only if both periods have defined endpoints. Overlapping subsystem work means \(B\) is not automatically a project delay. If no comparable EPS exists, time a small EPS design and change exercise to establish a local baseline.

    \item \textbf{Local target and forecast.} Have the advisor assign an EPS cost allowance \(C_{\mathrm{allow,EPS}}\) within the proposed mission funding plan and an allowable elapsed time \(T_{\mathrm{allow,EPS}}\) from a defined EPS start to the validation milestone. Estimate cost \(\widehat C_{\mathrm{EPS}}\) and elapsed time \(\widehat T_{\mathrm{EPS}}\) on the same basis from the ÉTS baseline and documented changes in payload loads, orbit, solar and battery choices, make or buy decisions, team capacity, and required tests. Spacecraft size alone cannot justify scaling cost or time. State the forecast gaps as \(\Delta C=\widehat C_{\mathrm{EPS}}-C_{\mathrm{allow,EPS}}\) and \(\Delta T=\widehat T_{\mathrm{EPS}}-T_{\mathrm{allow,EPS}}\). Only positive, supported gaps justify a claim that current performance is insufficient for the proposed EPS allocation or milestone.
\end{enumerate}

\paragraph{Provisional problem statement.}
\begin{quote}
For its proposed 3U CUBICS mission, ÉTS must design and validate an EPS that supports the required operating modes within an EPS allowance of [\(C_{\mathrm{allow,EPS}}\)] and [\(T_{\mathrm{allow,EPS}}\)] months from the defined EPS start to validation. The team's prior [\(X\)U] EPS required [\(C_{\mathrm{past,EPS}}\)] and [\(T_{\mathrm{past,EPS}}\)] under [scope and conditions]. After adjustment for the proposed mission, its estimated EPS cost and validation time exceed those allocations by [\(\Delta C\)] and [\(\Delta T\)].
\end{quote}

The bracketed values are evidence gaps, not estimates. The final statement should identify the source and uncertainty of each number. If either forecast gap is not positive, remove that claim from the problem statement and retain the corresponding budget or deadline as a constraint.
